% Esquisse d'un programme (cote 119, copy 1): archive cote 119, site page 8 (PDF page 9), typescript page 6.
% typescript. Draft compared: EsquisseFr.pdf (text layer):208-272.
% AI-transcribed by Claude Opus 5.5 (claude-opus-5-5, adaptive thinking, effort high), Message Batches API, batch esq1 (msgbatch_01LNsogFyF3y6KSp9UbtJtBk), 2026-11-12; not yet checked by eye.
\scanpage{119/8}
\mspage{6}

projectives avec quatre points marqués, soit pour des courbes ellip-%
tiques (!)) n'ait jamais été bien explicité, par exemple la relation
entre le cas de genre 0 avec la géométrie de l'octaèdre, et celle du
tétraèdre. A fortiori les multiplicités modulaires $M_{0,5}$ (pour les
droites projectives avec cinq points marqués) et $M_{1,2}$ (pour les
courbes de genre 1 avec deux points marqués), d'ailleurs quasiment
isomorphes entre elles, semblent-elles terre vierge - les groupes
de tresses ne vont pas nous éclairer à leur sujet ! J'ai commencé à
regarder $M_{0,5}$ à des moments perdus, c'est un véritable joyau, d'une
géométrie très riche étroitement liée à celle de l'icosaèdre.

L'intérêt a priori d'une connaissance complète des deux
premiers étages de la tour (savoir, les cas où la dimension modulaire
$N = 3g-3+\nu$ est $\leq 2$) réside dans ce principe, que \uline{la tour entière se
reconstitue à partir des deux premiers étages}, en ce sens que via
l'opération fondamentale de ``recollement'', l'étage 1 fournit un sys-%
tème complet de générateurs, et l'étage 2 un système complet de re-%
lations. Il y a une analogie frappante, et j'en suis persuadé, pas
seulement formelle, entre ce principe, et le principe analogue de
Demazure pour la structure des groupes algébriques réductifs, si on
remplace le terme ``étage'' ou ``dimension modulaire'' par ``rang semi-%
simple du groupe réductif''. Le lien devient plus frappant encore,
si on se rappelle que le groupe de Teichmüller $T_{1,1}$ (dans le contexte
discret transcendant maintenant, et non dans le contexte algébrique
profini, où on trouve les complétions profinies des premiers) n'est
autre que $S\ell(2,\hand{\mathbb{Z}})$, i.e. le groupe des points entiers du schéma en
groupes simple de rang 1 ``absolu'' $S\ell(2)_{\hand{\mathbb{Z}}}$. Ainsi, \uline{la pierre de cons-%
truction fondamentale pour la tour de Teichmüller}, \uline{est essentiellement
la même que celle pour ``la tour'' des groupes réductifs de tous rangs} -
un groupe d'ailleurs dont on peut dire sans doute qu'il est présent
dans toutes les disciplines essentielles des mathématiques.

Ce principe de construction de la tour de Teichmüller n'est
pas démontré à l'heure actuelle - mais je n'ai aucun doute qu'il ne
soit valable. Il résulterait (via une théorie de dévissage des struc-%
tures stratifiées - en l'occurence les $M^{\wedge}_{g,\nu}$ - qui resterait à écrire,
cf. par. 5) d'une propriété extrêmement plausible des multiplicités

\margin{8}
